% !TEX program = xelatex
% 编译方式：在本目录执行两遍 xelatex（详见 docs/_manual_common/README.md）。
% 章节源文件位于 chapters/，由本文件 \input 引入。
% 源码依据（行号以 2026-08-17 检出为准）：
%   include/hacdcpf/resilience/{resilience_assessment,certified_restoration,resilience_dynamic_certification}.hpp
%   src/resilience/{resilience_assessment,resilience_restoration_mip,resilience_stage_milp,certified_restoration}.cpp
\documentclass[UTF8,fontset=fandol,a4paper,11pt]{ctexart}
\usepackage{../../_manual_common/hysim_manual}

\title{\textbf{交直流混合高弹性能源电力系统仿真平台}\\[0.5em]
       {\Large 弹性恢复技术手册}\\[0.3em]
       {\large 数学模型、恢复算法、认证机制与接口参数}}
\author{HySim-XJTU-HRPES（西安交通大学 HRPES 团队）}
\date{\today}

\begin{document}
\maketitle
\begin{abstract}
\noindent 本手册依据 HySim-XJTU-HRPES 平台弹性恢复模块
（\srcpath{src/resilience/} 与 \srcpath{include/hacdcpf/resilience/}）整理，覆盖三种可互换的
配电弹性恢复模型——启发式时序重构、严格多时段 LinDistFlow 恢复 MILP、RA 式分阶段灾害
MILP——以及移动储能（MESS）时空路由、有限动作认证恢复（主问题 MIP + 信息物理/MESS
可执行性 + 多保真动态 DAE 证书）。每个方法给出功能定位、源码等价数学实现、全部选项与结果字段的单位制和
默认值、验证依据及已知局限。文档风格与《潮流计算技术手册》《最优潮流技术手册》一致。

\smallskip
\noindent\textbf{核心口径}：恢复 MILP 为 LinDistFlow/有功口径，\fld{model\_scope}=
\texttt{"hybrid-acdc-restoration-milp"}；\fld{mip\_gap\_within\_tolerance} \emph{不等于}全局最优；
认证恢复的 L3 \fld{proof\_valid} 只在“已仿真情景/阈值/模型/时域”范围内成立，\textbf{非}全局
暂态稳定证明；L1/L2 常量为采样估计（\fld{constants\_are\_global\_bounds}=false）。全部近似与
回退写入 \fld{model\_scope}、\fld{model\_limitations} 与 \fld{ValidityFlags}。
\end{abstract}

\tableofcontents
\newpage

% 第 1 章：架构总览与恢复模型族
\section{架构总览与恢复模型族}\label{sec:res-overview}

\subsection{功能定位}
弹性恢复模块回答“扰动/灾害致多重故障后，如何在多时段内以开关重构、DER/储能调度与
移动储能路由最大限度恢复供电，以及所选恢复动作在动态上是否安全”。命名空间
\texttt{hacdcpf::analysis}。三种可互换恢复模型由枚举 \fld{DistributionResilienceModel}
（\srcpath{resilience\_assessment.hpp:16}）选择：
\begin{center}\footnotesize
\begin{tabular}{@{}L{4.6cm}L{5.0cm}L{4.4cm}@{}}
\toprule
模型 & 方法 & 入口\\
\midrule
\fld{HeuristicSequential} & 贪心时序重构（默认） & \srcpath{run\_distribution\_resilience\_assessment}\\
\fld{MultiPeriodMIPLinDistFlow} & 严格多时段恢复 MILP & \srcpath{run\_distribution\_resilience\_mip\_assessment}\\
\fld{RAStyleStageMILP} & 分阶段灾害隔离/重构 MILP & \srcpath{run\_distribution\_resilience\_stage\_milp\_assessment}\\
\bottomrule
\end{tabular}
\end{center}
入口行号：\srcpath{run\_distribution\_resilience\_assessment}（:468）、
\srcpath{run\_distribution\_resilience\_mip\_assessment}（:458）、
\srcpath{run\_distribution\_resilience\_stage\_milp\_assessment}（:463）、
\srcpath{apply\_distribution\_resilience\_demo\_data}（:452）。未经本轮逐式核验的认证恢复推导
不纳入主手册数学章节。

\subsection{时序数据流}
\begin{figure}[H]\centering
\begin{tikzpicture}[font=\footnotesize,
  box/.style={draw,rounded corners,minimum height=0.9cm,align=center,inner sep=3pt},
  arr/.style={->,thick}]
  \node[box] (fault) at (0,0) {故障序列\\\fld{DistributionResilienceFault}};
  \node[box] (step) at (3.7,0) {逐时步\\负荷/RES 曲线};
  \node[box] (rec) at (7.0,0) {重构/调度\\贪心 或 MILP};
  \node[box] (mess) at (10.4,0) {MESS 路由\\时空+SOC};
  \node[box] (idx) at (3.7,-2.1) {弹性指标\\恢复比·加权失供};
  \node[box] (cert) at (7.6,-2.1) {动态认证\\L1/L2/L3 证书（可选）};
  \draw[arr] (fault)--(step); \draw[arr] (step)--(rec); \draw[arr] (rec)--(mess);
  \draw[arr] (rec.south) |- (idx.west); \draw[arr] (idx)--(cert);
\end{tikzpicture}
\caption{弹性恢复时序数据流：故障序列$\to$逐时步$\to$重构/调度$\to$MESS 路由$\to$指标（$\to$动态认证）}
\end{figure}

\subsection{弹性指标}
聚合结果 \fld{DistributionResilienceResult}（:413）给出恢复度量。设时步 $t$ 需求 $D_t$、
供电 $P_t^{srv}$、恢复比 $\gamma_t=P_t^{srv}/D_t$，则总量与面积型弹性指数
\begin{equation}
\text{加权失供}=\sum_t \Big(\sum_{i}w_i\,s_{i,t}\Big)\Delta t,\qquad
\fld{resilience\_index}\ \sim\ \frac{\sum_t P_t^{srv}\Delta t}{\sum_t D_t\Delta t},
\end{equation}
对应 \fld{weighted\_unserved\_mwh}、\fld{total\_served\_mwh}/\fld{total\_demand\_mwh}、
\fld{resilience\_index}、\fld{avg\_restoration\_ratio}、\fld{peak\_shed\_mw}，切负荷按优先级
四级分解 \fld{shed\_by\_priority}=[Critical, High, Medium, Low]。

\subsection{诚实口径}
恢复 MILP 的能力由 \fld{DistributionResilienceModelStats::ValidityFlags}（:88）逐项声明。
默认 \fld{model\_scope}=\texttt{"hybrid-acdc-restoration-milp"}；典型缺省下
\fld{dc\_branch\_flow\_limits\_enforced}=false、\fld{converter\_transfer\_limits\_enforced}=false、
\fld{converter\_losses\_modelled}=false、\fld{reactive\_power\_modelled}=false、
\fld{protection\_logic\_modelled}=false、\fld{transient\_limits\_modelled}=false，而
\fld{radial\_topology\_enforced}=true。\fld{mip\_gap\_within\_tolerance} 只表示与已知界的相对
间隙在容差内，\textbf{不}等于全局最优（准确间隙见 \fld{model\_stats.mip\_gap}）。

% ── 通用理论层（修订版 2 新增，工程参考级）──
\section{韧性的理论基础与问题分层}\label{sec:res-theory}

\begin{theorynote}
本章给出韧性曲线、随机风险、恢复优化与动态资格的通用数学框架，不是源码逐行转写。当前实现的
精确变量、约束、截断和回退分别在第~\ref{sec:res-source-model}--\ref{sec:res-dynamic}~章说明；
本章超出实现的理论形式均以未实现扩展框标识。
\end{theorynote}

\subsection{可靠性与韧性不是同一统计量}
可靠性通常以平稳或年度暴露下的失效频率、持续时间和期望失供电量刻画；韧性研究的是给定高影响
事件下系统性能从扰动前、退化、应急稳定到恢复的完整轨迹。令 $Q(t)\in[0,1]$ 为归一化服务
水平，$t_0$、$t_d$、$t_r$ 分别为扰动开始、最低服务点和研究终点，则面积型指标为
\begin{equation}
 R_Q=\frac{1}{t_r-t_0}\int_{t_0}^{t_r}Q(t)\,\mathrm dt,
 \qquad
 L_Q=\int_{t_0}^{t_r}[1-Q(t)]\,\mathrm dt .
\end{equation}
平台离散实现取 $Q_t=P_t^{srv}/P_t^{dem}$，因此
\begin{equation}
 EENS=\sum_t\sum_i s_{it}\Delta t,
 \qquad
 R_E=\frac{\sum_t P_t^{srv}\Delta t}
              {\sum_t P_t^{dem}\Delta t}.
\end{equation}
$R_E$ 对应 \fld{resilience\_index}；它是需求加权面积比，不等同于
$|T|^{-1}\sum_tQ_t$（\fld{avg\_restoration\_ratio}）。两者只有在各时步需求相同才相等。

\subsection{从灾害到可执行恢复的因果链}
一个可审计研究必须保留下列条件链，而不能从风速直接跳到“恢复成功”：
\begin{equation}
 \xi\longrightarrow h_{k,t}(\xi)\longrightarrow
 p_{k,t}^{fail}\longrightarrow a_{k,t}\longrightarrow
 (z,u,p,E,x)\longrightarrow Q_t\longrightarrow
 \mathcal C_{DAE}.
\end{equation}
其中 $\xi$ 是情景种子，$h$ 是灾害暴露，$a$ 是元件可用状态，$z/u$ 是支路/母线状态，
$p$ 是功率，$E$ 是储能，$x$ 是 MESS 位置，$\mathcal C_{DAE}$ 是动作的动态证书。
本模块只在每一箭头都有结构化数据时给下游结果授信；缺少保护绑定、零负荷拾取模型或动态初始化时，
返回 \fld{Unresolved}/\fld{Failed}，而不是补造数值。

\subsection{确定性、随机与风险口径}
给定单一情景 $\omega$ 的运行损失记为
\begin{equation}
 L(\omega)=\sum_{t,i}w_i s_{it}(\omega)\Delta t.
\end{equation}
对 $N$ 个独立、同分布样本，经验均值、经验分位数及条件尾期望为
\begin{align}
 \widehat\mu_N&=N^{-1}\sum_{n=1}^{N}L_n,\\
 \widehat{VaR}_{\alpha}&=L_{(\lceil\alpha N\rceil)},\\
 \widehat{CVaR}_{\alpha}&=\frac{1}{|\mathcal T_\alpha|}
   \sum_{n\in\mathcal T_\alpha}L_n,
 \quad \mathcal T_\alpha=\{n:L_n\ge\widehat{VaR}_\alpha\}.
\end{align}
复核案例采用配对种子比较基线与干预，降低灾害实现差异带来的方差；它是事件条件经验风险，
没有事件年发生率，因而不得解释为年度 EENS。

\begin{gapnote}
恢复模块没有概率测度校准、置信区间、重要抽样、相关随机场或分布鲁棒风险优化。经验均值/VaR/CVaR
由复核工具后处理，不是 \fld{DistributionResilienceResult} 的通用求解目标。
\end{gapnote}

\subsection{优化层与物理层}
恢复决策是混合离散问题。严格入口求解
\begin{equation}
 \min_{z,u,p,E,x,s}\quad
 \sum_{t,i}\pi_i s_{it}\Delta t+C_{sw}+C_{travel}+C_{dispatch}
\end{equation}
并受拓扑、径向性、有功平衡、电压包络、资源容量和跨时段状态约束。这里的
``严格''指模型按所列约束构成并交给 MIP 后端，不表示模型等价于非线性三相 AC/DC 物理。
动态层随后把离散动作映射为 DAE 事件，检查
\begin{equation}
 f_{min}\ge\underline f,\qquad
 |\dot f|_{max}\le\overline r_f,\qquad
 V_{min}\ge\underline V,
\end{equation}
以及换流器电流观测门。优化可行、MIP gap 合格和动态安全是三个独立命题。

\subsection{四类证明义务}
\begin{enumerate}
 \item \textbf{数据义务}：稳定 ID、AC/DC 域、单位、时标和情景种子可追溯；
 \item \textbf{优化义务}：报告可行性、目标、界、gap、规模、后端和时限；
 \item \textbf{物理义务}：用潮流或 DAE 对选定状态/动作复核，并报告残差与阈值；
 \item \textbf{统计义务}：风险样本数、配对规则、收敛准则和未收敛状态完整披露。
\end{enumerate}
只有四项均满足，案例才可进入“研究结论”；任何一项失败只能进入诊断结果。

\subsection{文献坐标与实现边界}
恢复曲线采用 Bruneau 等提出的性能轨迹思想；配电系统多时段恢复与微网成岛通常使用
混合整数网络流或 LinDistFlow；径向性以单商品流/虚拟根约束表达；动态认证使用轨迹残差、
采样雅可比与情景阈值检查。上述理论是理解框架。平台当前实现的唯一精确契约仍是
\srcpath{src/resilience/} 中的方程装配，故本手册后续各章都区分“理论可用形式”和“当前已实现形式”。

\subsection{与实现的对应}
\begin{center}\small
\begin{tabular}{@{}L{5.0cm}L{9.0cm}@{}}
\toprule
理论条目 & 实现状态\\\midrule
离散服务曲线、EENS 与面积型韧性指数 &
\implfull{src/resilience/resilience\_assessment.cpp:run\_distribution\_resilience\_assessment}\\
多时段 hybrid AC/DC 恢复 MIP &
\implfull{src/resilience/resilience\_restoration\_mip.cpp:build\_mip\_skeleton}\\
分阶段隔离与灾后重构 &
\implfull{src/resilience/resilience\_stage\_milp.cpp:run\_distribution\_resilience\_stage\_milp\_assessment}\\
轨迹残差与情景 DAE 阈值认证 &
\implpart{L3 为情景/模型/阈值/时域特定认证；L1/L2 常量仅为采样估计}\\
经验均值、VaR、CVaR 与配对干预 &
\implpart{仅复核工具 tools/distribution\_resilience\_review\_case.cpp 实现；非公共恢复求解器目标}\\
相关随机场、置信区间和分布鲁棒恢复 & \implnone\\
动态反馈自动加入恢复 MIP 并迭代至闭环 & \implnone\\
\bottomrule
\end{tabular}
\end{center}

% ── 实现转写层（源码等价）──
% 第 2 章：启发式时序重构
\section{启发式时序重构}\label{sec:res-heuristic}

\subsection{功能定位}
默认模型 \fld{HeuristicSequential}（\srcpath{run\_distribution\_resilience\_assessment}，
\srcpath{resilience\_assessment.hpp:468}）在 \fld{horizon\_hours} 个时步上逐步演进：施加/修复
故障、隔离故障区、贪心闭合联络开关以重新带电失电负荷、调度储能与 DER，并统计逐步供电
与切负荷。它快速但不保证全局最优，适合大规模快速评估与作为 MILP 的对照基线。

\subsection{逐步流程}
第 $t$ 步（步长 \fld{time\_step\_hr}）：
\begin{enumerate}
  \item \textbf{负荷/资源缩放}：以 \fld{load\_profile} 与 \fld{renewable\_profile}
        （或 \fld{pv\_profile}/\fld{wind\_profile}，缺省用内置 24h 曲线）得 $D_{i,t}$、$P^{res}_t$；
  \item \textbf{故障演进}：更新在运故障集（\fld{outage\_start\_hr}$\to$\fld{repair\_duration\_hr}）；
  \item \textbf{隔离}：以开关/断路器隔离故障区，得到带电子图；
  \item \textbf{贪心重构}：按最多 \fld{max\_reconfig\_iterations} 次迭代，逐次闭合能连通失电
        负荷且保持辐射的联络开关（受 \fld{allow\_reconfiguration} 控制）；
  \item \textbf{调度}：储能/构网 VSC/DER 在带电岛内出力（\fld{enable\_storage\_dispatch}、
        \fld{enable\_islanding}）；
  \item \textbf{MESS}：以电气图作运输代理（\fld{use\_electrical\_graph\_as\_transport\_proxy}）
        或显式 \fld{transport\_edges}，按 \fld{mess\_travel\_speed\_kmph} 路由；
  \item \textbf{统计}：得 $P_t^{srv}$、$s_{i,t}$、\fld{restoration\_ratio}、拓扑桥/割点等。
\end{enumerate}

\subsection{每步产物}
逐步结果 \fld{DistributionResilienceStepResult}（:332）记录带电/失电、开关动作
（\fld{switch\_open\_actions}/\fld{switch\_close\_actions}）、故障区母线、MESS 状态
\fld{mess\_states}、按优先级分解的 \fld{shed\_by\_priority}，以及拓扑韧性诊断
\fld{cut\_vertex\_bus\_ids}（割点，优先派修）与 \fld{bridge\_branch\_ids}（桥，恢复后重连最大子树）。
可选 \fld{run\_power\_flow}=true 时逐步跑完整潮流得电压/支路潮流（\fld{bus\_voltages}/\fld{branch\_flows}）。

\section{源码等价数学实现}\label{sec:res-source-model}

本章只覆盖可由当前实现逐行重建的启发式供电计算和严格多时段 MIP。认证恢复的动态证书由
\srcpath{src/resilience/resilience\_dynamic\_certification.inc} 实现；在未完成逐式审计前，本手册
不复述其理论误差界。

\subsection{启发式岛内供电分配}

若岛含电网源，\srcpath{evaluate\_component} 直接令供电等于总需求，切负荷为零；否则先把岛内
发电机、静止发电机、可再生和光伏容量相加为 $P_{avail}$。固定储能和 MESS 只放电、不充电，
对每个资源依次计算
\begin{align}
 P_{deliver}&=\min\!\left(P_{max},
 \max\!\left(0,\frac{(E-E_{min})\eta_{dis}}{\Delta t}\right)\right),\\
 P_{dispatch}&=\min(P_{remaining},P_{deliver}),
\end{align}
MESS 还必须已到外部可用时刻、非在途且位于该岛。然后
\begin{equation}
 P_{served}=\min(P_{demand},P_{avail}),\qquad
 P_{shed}=\max(0,P_{demand}-P_{served}).
\end{equation}
负荷按 \fld{importance} 降序、需求降序分配供电；每条负荷的加权切负荷为
$P_{shed,i}\texttt{importance}_i$。重要度阈值按 $4-10^{-9}$、$2.5-10^{-9}$、
$1.5-10^{-9}$ 划入关键、高、中、低四档。依据为
\srcpath{src/resilience/resilience\_assessment.cpp:StepContext}。

资源状态更新只有放电项：
\begin{equation}
 E\leftarrow\max(E_{min},E-P_{dispatch}\Delta t/\eta_{dis}).
\end{equation}
代码没有启发式充电变量，旧手册中的
$E_{t+1}=E_t+(\eta_{ch}P^{ch}-P^{dis}/\eta_{dis})\Delta t$ 不是该入口行为。
依据为 \srcpath{src/resilience/resilience\_assessment.cpp:apply\_dispatch\_to\_storage}。

\subsection{启发式 MESS 路由}

候选目标只来自“无电网源且切负荷大于 $10^{-9}$”的岛，其缺额使用
\fld{weighted\_shed\_mw}。每个空闲 MESS 对 Dijkstra 最短路执行以下资格判断：路由至少两个
节点，距离有限且大于 $10^{-9}$，累计距离不超过上限加 $10^{-9}$，并满足
$E-d e_{travel}>E_{min}+10^{-9}$。速度为
$\max(\fld{mess\_travel\_speed\_kmph},1)$，选择分数为
\begin{equation}
 score=\frac{\texttt{target.deficit}}
 {1+(d/v)/\max(\Delta t,10^{-9})}.
\end{equation}
仅在分数比当前最佳值大 $10^{-9}$ 时替换目标。每完成一段距离 $d_l$，状态执行
$E\leftarrow\max(E_{min},E-d_le_{travel})$；行程时间逐段按 $d_l/v$ 扣减。
依据为 \srcpath{src/resilience/resilience\_assessment.cpp:mark\_mess\_deployed}。

\subsection{启发式指标聚合}

每步汇总需求、供电、切负荷与加权切负荷，并计算
$r_t=P_{served,t}/P_{demand,t}$；当需求不大于 $10^{-9}$ 时 $r_t=1$。结果执行
\begin{align}
 E_{demand}&=\sum_tP_{demand,t}\Delta t,&
 E_{served}&=\sum_tP_{served,t}\Delta t,\\
 E_{shed}&=\sum_tP_{shed,t}\Delta t,&
 E_{weighted}&=\sum_tP_{weighted\ shed,t}\Delta t,\\
 RI&=\begin{cases}E_{served}/E_{demand},&E_{demand}>10^{-9},\\1,&\text{否则},\end{cases}&
 \overline r&=|T|^{-1}\sum_tr_t.
\end{align}
兼容字段 \fld{total\_ens\_mwh} 被直接设为 \fld{total\_shed\_mwh}，
\fld{restoration\_time\_hr} 被设为最后一步的 hour，而不是首次完全恢复时刻。
依据为 \srcpath{src/resilience/resilience\_assessment.cpp:run\_distribution\_resilience\_assessment}。

\subsection{严格 MIP 的实际变量}

\srcpath{build\_mip\_skeleton} 每个支路和时步创建：状态 $z_{bt}$、开/合动作
$y^{on}_{bt},y^{off}_{bt}$、有功 $p_{bt}$、商品流 $f_{bt}$；每个母线创建平方电压 $v_{it}$、
带电状态 $u_{it}$、根 $r_{it}$、基础根 $r^{base}_{it}$、根商品流、基础源出力、切负荷、
已用新能源和可调发电。固定储能只创建放电、能量和构网根；MESS 创建能量、位置、构网根、
放电和时空弧变量。当前文件没有无功 $Q_{bt}$ 或充电变量。
精确索引见 \srcpath{src/resilience/resilience\_restoration\_mip.cpp:IndexMaps} 和
\srcpath{src/resilience/resilience\_restoration\_mip.cpp:add\_var（lambda）}。

\subsection{严格 MIP 的目标系数}

线性目标只有以下非零系数：
\begin{align}
 c(y^{on}_{bt})=c(y^{off}_{bt})&=\fld{switching\_cost},\\
 c(s_{it})&=\pi_i\Delta t,\\
 c(P^{MESS}_{mit})&=10^{-4}\Delta t,\\
 c(a_{ma})&=\begin{cases}0,&\text{停留弧},\\
 \fld{mess\_travel\_cost\_per\_km}d_a+10^{-4},&\text{移动弧}.
 \end{cases}
\end{align}
$\pi_i$ 按四级优先级返回对应 \fld{shed\_penalty\_*}。当前
\fld{renewable\_curtailment\_cost} 未进入本文件的目标系数，不能写入现行目标。
依据为 \srcpath{src/resilience/resilience\_restoration\_mip.cpp:resilience\_mip\_solver\_name}、
\srcpath{src/resilience/resilience\_restoration\_mip.cpp:penalty\_from\_priority} 和
\srcpath{src/resilience/resilience\_restoration\_mip.cpp:add\_var（lambda）}。

\subsection{支路动作、流限与电压行}

初始时步与后续时步分别创建
\begin{align}
 z_{b0}-y^{on}_{b0}+y^{off}_{b0}&=z_b^{initial},\\
 z_{bt}-z_{b,t-1}-y^{on}_{bt}+y^{off}_{bt}&=0,\quad t>0,\\
 y^{on}_{bt}+y^{off}_{bt}&\le1,\\
 -\bar P_bz_{bt}\le p_{bt}&\le\bar P_bz_{bt},\\
 -n_bz_{bt}\le f_{bt}&\le n_bz_{bt},\\
 z_{bt}&\le u_{from,t},\quad z_{bt}\le u_{to,t}.
\end{align}
只对 \fld{enforce\_voltage\_drop}=true 的边创建
\begin{equation}
 |v_{to,t}-v_{from,t}+2r_bp_{bt}|\le
 \fld{big\_m\_voltage}(1-z_{bt}).
\end{equation}
没有 $x_bQ_{bt}$ 项。依据为
\srcpath{src/resilience/resilience\_restoration\_mip.cpp:add\_le（lambda）}。

\subsection{母线平衡与森林行}

每时步只创建一次全网森林基数等式：
\begin{equation}
 \sum_bz_{bt}+\sum_ir_{it}-\sum_iu_{it}=0.
\end{equation}
母线有功行的源码符号为
\begin{equation}
 P^{base}_{it}+P^{gen}_{it}+P^{ren}_{it}+s_{it}
 +\sum_{b:to=i}p_{bt}-\sum_{b:from=i}p_{bt}
 +\sum_{s@i}P^{fs}_{st}+\sum_mP^{MESS}_{mit}=D_{it}.
\end{equation}
商品流行是
$F^{root}_{it}-u_{it}+\sum_{to=i}f_{bt}-\sum_{from=i}f_{bt}=0$。
另有 $u_i-r_i-\sum_{b\ni i}z_b\le0$、$r_i\le u_i$、基础源/发电/新能源的带电上界，
以及 $s_{it}\ge D_{it}(1-u_{it})$。根电压通过两条 big-M 行固定到 1，带电电压下界为
$v_{it}\ge v_{min}^{2}u_{it}$；变量自身上界为 $v_{max}^{2}$，代码未创建
$v_{it}\le v_{max}^{2}u_{it}$。依据为
\srcpath{src/resilience/resilience\_restoration\_mip.cpp:add\_le（lambda）}。

\subsection{储能和 MESS 能量行}

固定储能初值固定为 $E_{s0}=E_s^{init}$，并创建
\begin{align}
 P^{fs}_{st}&\le P_s^{max}u_{i(s),t},\\
 (\Delta t/\max(\eta_{dis},10^{-6}))P^{fs}_{st}-E_{st}&\le-E_s^{min},\\
 E_{st}-E_{s,t-1}+(\Delta t/\max(\eta_{dis},10^{-6}))P^{fs}_{s,t-1}&=0.
\end{align}
最后一式仅对 $t>0$ 创建，因此当前时步 $t$ 的放电从下一时步能量中扣除。

MESS 同样固定 $E_{m0}$，每步位置满足 $\sum_ix_{mit}=1$，放电满足位置上界，能量可用行是
\begin{equation}
 -E_{mt}+\frac{\Delta t}{\max(\eta_{dis},10^{-6})}\sum_iP^{MESS}_{mit}
 \le-E_m^{min}.
\end{equation}
$t>0$ 的递推还将上一时步出发的移动弧能耗加入：
\begin{equation}
 E_{mt}-E_{m,t-1}+\frac{\Delta t}{\max(\eta_{dis},10^{-6})}
 \sum_iP^{MESS}_{mi,t-1}+\sum_{a:\,depart=t-1}e_a a_{ma}=0.
\end{equation}
位置转移用“每个位置恰选一条出发弧”和“每个位置由到达弧确定”两组等式；见
\srcpath{src/resilience/resilience\_restoration\_mip.cpp:add\_le（lambda）}。

\subsection{求解结果真实性}

返回可行不仅依赖后端 success，还要求解向量长度正确，变量界、不等式、等式、整数性最大违反
均不超过 $10^{-6}$；若所有目标系数不小于 $-10^{-9}$，重算目标还必须不小于 $-10^{-6}$。
\fld{mip\_gap\_within\_tolerance} 使用
$\texttt{success}\land gap\le\fld{mip\_gap}+10^{-9}$。依据为
\srcpath{src/resilience/resilience\_restoration\_mip.cpp:solve\_branch\_and\_cut\_isolated（lambda）}。

\subsection{覆盖边界}

RA 式分阶段入口由 \srcpath{src/resilience/resilience\_stage\_milp.cpp} 单独构建，不能假定与本章
严格多时段矩阵完全相同。其隔离、重构和 MESS 回退语义由
\fld{model\_scope}、\fld{mess\_dispatch\_model} 及有效性标志报告；本章不以严格 MIP 公式代替它。

% 第 5 章：RA 式分阶段灾害 MILP
\section{RA 式分阶段灾害 MILP}\label{sec:res-stage}

\subsection{功能定位}
\fld{RAStyleStageMILP}（\srcpath{run\_distribution\_resilience\_stage\_milp\_assessment}，
\srcpath{include/hacdcpf/resilience/resilience\_assessment.hpp:run\_distribution\_resilience\_stage\_milp\_assessment}；由 \fld{use\_ra\_style\_stage\_milp} 或首选开关
\fld{enable\_disaster\_stages} 触发）以“灾害隔离 + 灾后重构”两阶段、开关约束的拓扑 MILP
求解，阶段语义见 \fld{DistributionDisasterStage}（:DistributionDisasterStage）：
\texttt{Normal} $\to$ \texttt{DisasterIsolation} $\to$ \texttt{DisasterPostFaultReconfig}
$\to$ \texttt{PostDisasterRepair}。

\subsection{分阶段模型}
\begin{itemize}
  \item \textbf{阶段 1 灾害隔离}：以开关隔离受灾支路
        （\fld{use\_switch\_based\_fault\_isolation}=true，\fld{allow\_stage1\_open\_switches}），
        非故障支路的操作须有开关（\fld{require\_switch\_for\_nonfault\_branch\_operation}），
        默认禁止无开关直接操作支路（\fld{allow\_branch\_operation\_without\_switch}=false）；
  \item \textbf{阶段 2 灾后重构}：在 \fld{post\_fault\_reconfig\_window\_hr}（默认 2h）内闭合联络
        （\fld{allow\_stage2\_close\_ties}）以恢复负荷，可限 \fld{use\_remote\_switch\_only}；
  \item 每阶段使用有功输运、节点供需、带电与径向拓扑约束；不建立严格多时段 MIP 的 AC
        支路额定流限或电压包络，故相应 validity flags 为 false。
\end{itemize}
MESS 既可由专用严格 MILP（\fld{use\_strict\_mip\_for\_mess}）也可阶段内调度并在失败时回退
（\fld{fallback\_to\_stage\_mess\_dispatch}）。逐步结果记灾害阶段 \fld{disaster\_stage}、隔离性开断
\fld{isolation\_open\_ac\_branch\_ids}/\fld{isolation\_open\_dc\_branch\_ids} 与重构性合断，并区分
\fld{isolation\_switch\_actions} 与 \fld{reconfiguration\_switch\_actions}。

\subsection{与其他模型的关系}
三种模型互斥选择、共享同一逐步结果结构与弹性指标；启发式最快、多时段 MILP 全局协调、
RA 式分阶段贴合灾害隔离-重构的工程流程。三者的能力边界均由 \fld{ValidityFlags} 显式声明，
拓扑求解失败显式反映；MESS 子问题若允许回退则在 \fld{mess\_dispatch\_model} 与状态中明确标识。

% 第 7 章：选项与结果数据结构
\section{选项与结果数据结构}\label{sec:res-io}

\subsection{弹性研究选项 \texttt{DistributionResilienceOptions}}
\compmeta{DistributionResilienceOptions}{include/hacdcpf/resilience/resilience\_assessment.hpp:DistributionResilienceOptions}
\begin{paramtable}
\fld{model} & — & — & 枚举 & Heuristic & 恢复模型选择（3 选 1）\\
\fld{horizon\_hours} & $T$ & 小时 & $24\sim168$ & 48 & 研究时域\\
\fld{time\_step\_hr} & $\Delta t$ & 小时 & $0.5\sim1$ & 1.0 & 时步长\\
\fld{load\_scale\_factor} & — & — & $\ge1$ & 1.0 & 负荷加压系数\\
\fld{allow\_reconfiguration} & — & — & 布尔 & true & 允许开关重构\\
\fld{allow\_mess\_dispatch} & — & — & 布尔 & true & 允许 MESS 调度\\
\fld{max\_reconfig\_iterations} & — & — & 整数 & 50 & 贪心重构每步迭代上限\\
\fld{mess\_travel\_speed\_kmph} & $v$ & km/h & $20\sim60$ & 40 & MESS 行驶速度\\
\fld{default\_fault\_count} & — & — & 整数 & 2 & 自动生成故障数\\
\fld{default\_repair\_time\_hr} & $\tau_{rep}$ & 小时 & — & 6 & 默认修复时长\\
\fld{run\_power\_flow} & — & — & 布尔 & false & 逐步跑完整潮流校验\\
\fld{enable\_disaster\_stages} & — & — & 布尔 & false & 启用分阶段灾害隔离/重构\\
\fld{post\_fault\_reconfig\_window\_hr} & — & 小时 & — & 2.0 & 灾后重构窗口\\
\fld{use\_strict\_mip\_for\_mess} & — & — & 布尔 & true & MESS 专用严格 MILP\\
\end{paramtable}

\subsection{恢复 MILP 选项 \texttt{DistributionResilienceMIPOptions}}
\compmeta{DistributionResilienceMIPOptions}{include/hacdcpf/resilience/resilience\_assessment.hpp:DistributionResilienceMIPOptions}
\begin{paramtable}
\fld{v\_min\_pu}/\fld{v\_max\_pu} & $\underline v,\overline v$ & pu & $0.95/1.05$ & 0.95/1.05 & 电压包络\\
\fld{big\_m\_voltage} & $M$ & — & — & 4.0 & LinDistFlow 压降 big-M\\
\fld{default\_branch\_rate\_mva} & $\bar S_e$ & MVA & — & 10.0 & 缺省支路容量\\
\fld{shed\_penalty\_critical} & $\pi$ & — & — & $10^6$ & 关键负荷切负荷罚权\\
\fld{shed\_penalty\_high} & $\pi$ & — & — & $10^4$ & 高优先级罚权\\
\fld{shed\_penalty\_medium} & $\pi$ & — & — & $10^2$ & 中优先级罚权\\
\fld{shed\_penalty\_low} & $\pi$ & — & — & 1 & 低优先级罚权\\
\fld{switching\_cost} & $c_{sw}$ & — & — & 5.0 & 单次开关操作代价\\
\fld{mess\_travel\_cost\_per\_km} & $c_{mess}$ & — & — & 0.1 & MESS 行驶代价\\
\fld{renewable\_curtailment\_cost} & $c_{curt}$ & — & — & 10.0 & 弃风光代价\\
\fld{mip\_gap} & — & — & $0\sim0.05$ & 0.02 & MILP 相对间隙\\
\fld{max\_time\_s} & — & 秒 & — & 120 & 求解时限\\
\fld{max\_nodes} & — & — & — & 50000 & B\&C 节点上限\\
\fld{solver} & — & — & 枚举 & Native & MILP 后端\\
\fld{enable\_cglp\_cuts} & — & — & 布尔 & false & CGLP 提升-投影割\\
\end{paramtable}

\subsection{动态证书选项 \texttt{MultiFidelityCertificateOptions}}
\compmeta{MultiFidelityCertificateOptions}{include/hacdcpf/resilience/certified\_restoration.hpp}
\begin{paramtable}
\fld{min\_frequency\_hz} & — & Hz & — & 49.0 & L3 频率下限阈值\\
\fld{max\_rocof\_hz\_s} & — & Hz/s & — & 1.0 & L3 RoCoF 上限\\
\fld{min\_ac\_voltage\_pu} & — & pu & — & 0.90 & L3 AC 电压下限\\
\fld{require\_converter\_current\_observation} & — & — & 布尔 & true & L3 换流器电流观测门\\
\fld{jacobian\_tube\_factor} & — & — & — & 1.25 & 雅可比管因子\\
\fld{dense\_jacobian\_dimension\_limit} & — & — & — & 400 & 稠密雅可比维度上限\\
\fld{level1\_step\_multiplier} & — & — & — & 4.0 & L1 步长倍数\\
\fld{level2\_step\_multiplier} & — & — & — & 2.0 & L2 步长倍数\\
\fld{max\_initial\_dynamic\_residual} & — & — & — & $10^{-6}$ & 初始动态残差上限\\
\end{paramtable}

\subsection{结果结构摘要}
\begin{itemize}
  \item \fld{DistributionResilienceResult}（:CertifiedRestorationCoordinator）：逐步 \fld{steps}、\fld{fault\_sequence}、
        \fld{total\_served\_mwh}/\fld{total\_shed\_mwh}/\fld{weighted\_unserved\_mwh}、
        \fld{resilience\_index}、\fld{mess\_energy\_delivered\_mwh}、\fld{model\_stats}；
  \item \fld{DistributionResilienceStepResult}（:CertifiedRestorationCoordinator）：\fld{restoration\_ratio}、
        \fld{shed\_by\_priority}、开关动作、\fld{mess\_states}、桥/割点、可选潮流校验；
  \item \fld{DistributionResilienceModelStats}（:MultiFidelityCertificate）：MILP 规模/状态与 \fld{ValidityFlags}（20+ 标志）；
  \item \fld{CertifiedRestorationResult}：\fld{optimality\_proven\_over\_catalog}、\fld{no\_good\_cuts}、
        逐轮 \fld{iterations} 与 \fld{MultiFidelityCertificate}。
\end{itemize}

% 第 8 章：验证与边界局限
\section{验证与边界局限}\label{sec:res-verify}

\subsection{注册测试与复现}
\begin{center}\footnotesize
\begin{tabular}{@{}L{5.6cm}L{4.0cm}L{5.0cm}@{}}
\toprule
测试/目标 & 注册位置 & 覆盖\\
\midrule
\srcpath{test\_resilience\_assessment} & tests/CMakeLists.txt:272 & 三种恢复模型、MESS、指标、有效性\\
\srcpath{test\_transient\_dynamics} & tests/ & \srcpath{CertifiedRestorationCoordinator}、可执行性、MESS/信息物理阻断\\
\srcpath{cyber\_dynamic\_safe\_restoration\_small} & 可执行目标 & 认证恢复小案例复现\\
\srcpath{cyber\_dynamic\_safe\_restoration\_scale} & 可执行目标 & 认证恢复规模复现\\
\bottomrule
\end{tabular}
\end{center}
复现命令（\srcpath{docs/theory/certified\_restoration\_runtime.md}）：
\begin{quote}\footnotesize
\texttt{cmake --build build/macos-release --target}\\
\srcpath{cyber\_dynamic\_safe\_restoration\_small}\\
\srcpath{cyber\_dynamic\_safe\_restoration\_scale}
\end{quote}

\subsection{边界与局限（如实标注）}
\begin{itemize}
  \item \textbf{LinDistFlow/有功口径}：恢复 MILP \fld{model\_scope}=\texttt{"hybrid-acdc-restoration-milp"}，
        默认不强制 DC 支路流限、换流器传输限、换流器损耗与无功
        （\fld{ValidityFlags}，:ValidityFlags），强制辐射拓扑；
  \item \textbf{近似最优}：\fld{mip\_gap\_within\_tolerance}=true 只表示与已知界的间隙在容差内，
        \textbf{非}全局最优；准确间隙见 \fld{model\_stats.mip\_gap}；
  \item \textbf{认证 L1/L2}：采样雅可比“管”诊断，\fld{constants\_are\_global\_bounds}=false，
        须以验证包络替换采样量后方为解析剪枝证明；
  \item \textbf{认证 L3}：\fld{proof\_valid} 只在已仿真情景/阈值/模型/时域内成立，\textbf{非}全局
        暂态稳定证明；目录最优性仅覆盖有限目录内动作；
  \item \textbf{事件映射}：VSC/DC-DC 闭合、LCC 拓扑变化、独立 DC 断路器动作无精确动态事件，
        \fld{require\_supported\_transitions}=true 时标 \texttt{unresolved} 且不能置 \fld{proof\_valid}；
  \item \textbf{MESS}：行程动态不在亚秒 DAE 时间轴积分，到达/行程能量为主问题级前提；
  \item \textbf{桥接器}：仅重放拓扑驱动的带电，不回填完整 MIP 运行点，聚合出力不分配到单个
        动态器件（\fld{DistributionResilienceDAECertificateResult::limitations}）。
\end{itemize}

\subsection{加深与后续}
行号以 2026-08-17 检出为准，如与后续变更不符以
\srcpath{include/hacdcpf/resilience/} 与注册测试为准。数学与运行契约对照
\srcpath{docs/theory/certified\_restoration\_runtime.md}；网络重构建模见
\srcpath{docs/theory/network\_reconfiguration\_models.md}。


\section{跨模块工业级技术评价基线}

本章规定本手册所述模块进入工程算例、批量研究或生产服务前必须共同满足的评价基线。
具体模型公式、变量、约束和专用算法以正文相应章节为准；本章不扩大源码已经实现的范围。

\subsection{输入、输出、边界条件与数据结构审计}
输入审计应逐项检查有限性、单位、上下界、枚举值和引用完整性。系统对象必须先按所需层级完成
校验；AC 与 DC 母线采用域限定映射，组件稳定 \texttt{index}、向量位置和临时图索引不得混用。
输出审计必须记录单位、索引空间、模型范围、收敛或可行状态、后端、容差、迭代/节点计数和告警。
空系统、空时域、孤岛、重复 ID、零容量、退化约束与不可用可选后端均属于边界测试集。

\subsection{工程假设与数值稳定性分析}
模型使用的平衡、线性化、稳态、独立性、概率分布、时间离散或网络等值假设必须与应用场景匹配。
数值评价至少包含变量缩放、绝对/相对残差、可行性违反、条件数或枢轴质量、步长/阻尼、迭代停滞
和随机种子复现性。若采用正则化、平滑、回退、松弛或近似，应同时报告触发条件、参数值、对主指标
的敏感性以及是否改变模型口径；不得用“已返回结果”替代收敛或有效性证明。

\subsection{复杂度与资源分析}
令 $n_x$、$n_c$、$n_t$、$n_s$、$|V|$、$|E|$ 分别表示变量、约束、时段、场景、节点和边数。
报告应分别给出建模、求解、后处理的时间与内存。图遍历通常为 $O(|V|+|E|)$；逐时段/逐场景
流程至少线性增长为 $O(n_t n_s)$ 次子问题调用；稠密线性代数最坏可达 $O(n_x^3)$，稀疏方法则应
报告结构非零数、填充率和因子分解峰值内存；MILP/NLP 不给出虚假的多项式最坏界，而报告节点数、
间隙、时限和实测扩展曲线。复杂度结论必须基于固定硬件、固定后端和可复现算例族。

\subsection{异常工况处理}
非法输入应在进入求解器前返回结构化错误；不可行、无界、不收敛、时限、内存不足和后端缺失必须
相互区分。部分结果只有在对应 \fld{ValidityFlags}、\fld{model\_scope}、
\fld{model\_limitations} 或等价字段允许时才可使用。异常路径不得静默丢弃 AC/DC 资产、制造平衡
节点、把 NaN 改写为零或把近似解标为全局最优；系统原始对象应保持可恢复和可重试。

\subsection{与其他模块的接口关系}
接口链路遵循“工程语义模型 $\to$ 校验 $\to$ 投影/装配 $\to$ 求解或分析 $\to$ 结果回投影”。
跨模块传递时保留稳定 ID、单位、符号、时间基准和场景身份；消费潮流、OPF、动态或可靠性结果前，
先检查其收敛、覆盖范围和有效性标志。HTTP/JSON/Python 层只做语义保持适配，不重新解释求解结果。

\subsection{测试与验证建议}
最低验证矩阵包括：手算小例、标准算例、边界/异常输入、随机性质测试、算法或后端交叉校核、
序列化往返、稳定 ID 回投影、AC/DC 同号母线、重复运行确定性和规模扩展测试。数值算法还应加入
残差独立重算、雅可比有限差分或自动微分校核；近似模型应与高保真模型比较并给出误差分布，而非
只报告单个平均值。回归报告必须列明构建配置、算例版本、容差、通过/失败/跳过数及未验证范围。

\subsection{工业级评估指标}
\begin{itemize}
  \item \textbf{正确性}：守恒残差、约束最大违反、解析/基准误差、结果回投影一致性；
  \item \textbf{鲁棒性}：收敛率、不可行识别率、异常分类准确率、扰动敏感性与随机复现率；
  \item \textbf{性能}：端到端时延、吞吐量、峰值内存、并行效率与规模增长斜率；
  \item \textbf{可追溯性}：输入模式版本、稳定 ID 覆盖率、模型范围/限制字段完整率、告警可定位率；
  \item \textbf{工程有效性}：与标准、外部引擎或现场数据的偏差，以及误判/漏判对业务指标的影响。
\end{itemize}
指标应预先给出验收阈值和失败处置；未达到阈值时只能标记为实验性或受限适用。

\subsection{适用范围、局限性与后续改进}
本手册只评价当前源码覆盖的模型、后端和数据结构，不对未建模物理过程、未安装求解器或未执行的
外部验证作保证。后续改进应按“缺口 $\to$ 可量化指标 $\to$ 源码/测试变更 $\to$ 独立验证”闭环，
优先消除会造成错误结论的语义缺口，其次提升数值鲁棒性与性能，最后扩展模型范围。


\end{document}
